\section{The complete cubic graded model at the transported local component}
\label{sec:cubic-full-model}

The actual ring and relative presentation of AC-G49--G50
can be converted into the source's full graded model.
This requires checking its integral derivative field, not
only its dimensions.  We make that check explicitly.
All hypotheses of AC-G50, including \texttt{CohInp}, the
old \texttt{GD} conditions and the full \texttt{LemmaP},
remain in force.  The coefficient prime is in the range
of AC-G49-inert.  Put $E=\operatorname{Frac}(O)$ and
$\overline E=\overline{\mathbb Q}_\ell$ with the source's
specified coefficient embedding.

All source paths below have the pinned root
\path{Development/AutomorphicArithmetic/AutomorphicVarieties/}.
No new literal instance of the rational-character
\texttt{Dat} record is being asserted: the result verifies
every field of the corresponding \texttt{Model} definition
for the actual cubic ring and its actual basepoint.

\begin{lemma}[AC-G51-diagonal: integral derivatives at the cubic basepoint]
\label{AC-G51-diagonal}
Let $R$ be a commutative $O$-algebra, $x:R\to O$ an
$O$-algebra homomorphism, $Q$ an $O$-module, and
$\delta:R\to Q$ an $O$-linear derivation at $x$:
\[
 \delta(fg)=x(f)\delta(g)+x(g)\delta(f).
\]
Suppose $P\in\operatorname{GL}_3(R)$ and
$x(P)=\operatorname{diag}(d_0,d_1,d_2)$ with every
$d_i\in O^\times$.  Then
\[
 \begin{split}
 \delta((P^{-1})_{pq})
       &=-d_p^{-1}d_q^{-1}\delta(P_{pq}),\\
 \delta((\theta P)_{ij})
       &=-d_{j^{\rm rev}}^{-1}d_{i^{\rm rev}}^{-1}
                       \delta(P_{j^{\rm rev},i^{\rm rev}}),
       \qquad i^{\rm rev}=2-i .
 \end{split}                                                   \tag{G51.1}
\]
Consequently, at the cubic basepoint, the derivative of
each six-dimensional embedded entry difference is either
zero or an integral unit multiple of the derivative of
an actual three-dimensional entry of the same weight.
The multiplier depends only on the basepoint matrix,
the multiplier sign and the indices.  It is the same
for the global entry and its restricted local entry,
and is compatible with reduction modulo any nonzero
$a\in O$.
\end{lemma}
\begin{proof}
The derivation rule applied to $1\cdot1$ gives
$\delta(1)=0$, and $O$-linearity then gives
$\delta(O)=0$.  Apply $\delta$ entrywise to
$PP^{-1}=I$.  The $(p,q)$ entry of the resulting
identity is
\[
       d_p\delta((P^{-1})_{pq})
                    +d_q^{-1}\delta(P_{pq})=0,
\]
because $x(P)$ and $x(P^{-1})$ are diagonal.
Multiplication by the unit $d_p^{-1}$ proves the first
identity.  The source involution is
$\theta(P)=JP^{-t}J$, whose $(i,j)$ entry is
$(P^{-1})_{j^{\rm rev},i^{\rm rev}}$.  This proves the
second identity.

The six-dimensional embedding has blocks $P$, $\theta P$
and zero, with their positions determined by the multiplier
sign.  Its derivative is therefore respectively an actual
entry derivative, the second expression in (G51.1), or zero.
For the weights $w_i=1-i$,
\[
             w_{j^{\rm rev}}-w_{i^{\rm rev}}=w_i-w_j.
                                                               \tag{G51.2}
\]
Thus the reversed entry has the required weight.
At the cubic basepoint the three $d_i$ are
$q(g)^j\gamma(g),1,q(g)^j\gamma(g)^{-1}$, all units.
Subtracting a basepoint entry does not alter its derivative.
The displayed unit factors are algebraic expressions in
these actual basepoint entries, so restriction and reduction
preserve them.  This generalizes the explicitly verified
diagonal calculation in
\path{UnitaryThreeGroups/ExistsPointDerivationEmbEntryEqSmulMatEntry.lean},
\texttt{der\_inv} and
\texttt{exists\_pointDerivation\_emb\_entry\_eq\_smul\_mat\_entry}.
It does not require the special relation $d_0d_2=1$.
\end{proof}

\begin{theorem}[AC-G51: every field of the actual cubic model]
\label{AC-G51}
Let $S=R_\chi$ be the normalized ring of AC-G49 and
$x_0:S\to O$ its cubic basepoint.  With the notation
$a_i$ of AC-G50, set
\[
 \begin{split}
 r&=a_{-1}+2,\qquad s=a_{-2}+1,\\
 r'&=a_1+3,\qquad s'=a_2+1,\qquad m=a_0+3 .
 \end{split}                                                   \tag{G51.3}
\]
There is a graded $O$-algebra
$A=\bigoplus_{i\in\mathbb Z}\mathcal A_i$, an origin
$x:A\to O$, a homomorphism $q_A:A\to S$, and homogeneous
equations
\[
 \begin{aligned}
 F_1,\ldots,F_{r+1}&\in\mathcal A_{-1},&
 G_1,\ldots,G_s&\in\mathcal A_{-2},\\
 E_1,\ldots,E_{r'}&\in\mathcal A_1,&
 E'_1,\ldots,E'_{s'}&\in\mathcal A_2,\\
 Z_1,\ldots,Z_m&\in\mathcal A_0 .
 \end{aligned}                                                \tag{G51.4}
\]
with the following complete list of model properties.
Write $J$ for their generated ideal.

\begin{enumerate}
\item The degree-zero ring $\mathcal A_0$ has the source
\texttt{IsCO O} structure.  Finitely many elements of
weights $-2,-1,1,2$ generate $A$ over $\mathcal A_0$,
and $x$ vanishes in every nonzero weight.  Thus
\texttt{IsWRing} holds.

\item All equations vanish at $x$, and
$\ker q_A=J$, $x_0q_A=x$ and $r\ge1$.
There is a chart in the precise source sense
\texttt{SmoothPt}, over $E$, whose numbers of variables
of weights $(-2,-1,0,1,2)$ are
\[
                         (s,r,m,r'+1,s'+1).                   \tag{G51.5}
\]
This states the source chart condition; it does not
assert integral formal smoothness over $O$.

\item Every actual universal matrix entry of weight $i$
has a preimage in $\mathcal A_i$.  For every
$f\in\mathcal A_i$ and every $n$, $q_A(f)$ is congruent
modulo $\mathfrak m_S^n$ to an $O$-linear combination of
products of universal matrix entries of total weight $i$.

\item Retain the source condition
\[
 M2:\quad \operatorname{im}(\mathcal A_0\longrightarrow A/J)
                 \text{ is a finitely generated }O\text{-module}.
                                                               \tag{G51.6}
\]
Then every $\overline E$-valued point of $A$ killing $J$
is, after one weight-torus action, the restriction of an
actual local point of $S$.  Also every derivation at such
a local point which kills $J$ extends to an actual
dual-number point of $S$ with that first component.

\item Let $0\ne a\in O$.  Every derivation at $x_0$
with values in $O/a$ which vanishes on actual entries
of weights other than $-1$, and whose actual local
restriction has an integral weight-$-1$ derivation
annihilating $P_\chi$, lifts along $q_A$ to an integral
weight-$-1$ derivation at $x$.
\end{enumerate}

These are exactly the fields
\texttt{wring}, \texttt{Q}, \texttt{one\_le},
\texttt{smooth}, \texttt{ker\_q}, \texttt{base},
\texttt{x\_0\_q}, \texttt{ent}, \texttt{dense},
\texttt{pts}, \texttt{der} and \texttt{lift} of
\path{UnitaryThreeGroups/ConjStable.lean}, \texttt{Model},
with the actual cubic universal objects and basepoint.
The condition (G51.6) is an antecedent of the two
corresponding fields, not an additional conclusion.
\end{theorem}
\begin{proof}
\proofstep{Construct the ambient ring using only the local relations.}
Choose the actual local presentation $W_\ell$ and
relative presentation of AC-G50.  After renaming variables,
write
\[
 P=O[[X_s,Y_t]],\qquad
 \widetilde R=P/(I_\ell P).
                                                               \tag{G51.7}
\]
Only the local relations occur in $\widetilde R$.
The global relative equations and the eight homogeneous
normalization equations determine a surjection
$q:\widetilde R\twoheadrightarrow S$.
Choose a homogeneous lift to $P$ of each normalization
equation; such a lift exists by entry homogeneity and
the definition of the weight parts.  Let $g_u$ denote
their images together with the global equations in
$\widetilde R$.  Then
\[
                         \ker q=(g_u)_u.                     \tag{G51.8}
\]
Indeed the preimage of the normalized ideal in $P$ is
exactly the sum of the extended local ideal, the original
global relation ideal and these lifted equations.  This
proves both inclusions in (G51.8).
The extended local ideal is homogeneous; therefore
(G51.7) has the induced weighted presentation $W_A$.
Its coordinates map under $q$ to actual embedded entry
differences, including the local coordinates, by the
relative presentation compatibility in AC-G50.
This is the seed-independent construction
\texttt{mop\_cover} in
\path{UnitaryThreeGroups/ModelOfPres.lean}.

Let $A$ be the subalgebra of $\widetilde R$ consisting
of finite sums of homogeneous elements, with grading
$\mathcal A$, and put $q_A=q|_A$.
The cubic basepoint kills every coordinate, so
$x_0q$ is constant-term evaluation in (G51.7).
It therefore defines the origin $x$ on $A$.
All coordinate weights lie in $[-2,2]$.
The general weighted-presentation result
\path{CommutativeAlgebra/IsWRing.lean},
\texttt{isWRing}, gives exactly property 1 from this
bounded set of weights and the origin.  Its application
depends only on the weighted presentation, not its seed.

\proofstep{The equation ideal and the exact smooth-point dimensions.}
The homogeneous elements $g_u$ lie in $A$.
The homogeneous-ideal identity
\path{CommutativeAlgebra/HomogSpan.lean},
\texttt{homog\_span}, gives
\[
                  A\cap(g_u)\widetilde R=(g_u)A.
                                                               \tag{G51.9}
\]
Thus (G51.8) implies $\ker q_A=(g_u)A$.
All $g_u$ vanish at the origin because they vanish in
$S$ and $x=x_0q_A$.
There are no nonzero equation counts outside $[-2,2]$:
the extra-coordinate count there is zero and
\texttt{boxDef} is zero.  The exact counts in AC-G50
are precisely those in (G51.4) with (G51.3).
Grouping the finitely many $g_u$ by weight therefore
constructs (G51.4) and $J$, without choosing a new ideal.
This is the complete elementary enumeration in
\path{UnitaryThreeGroups/ExistsEqnsOfWeightedGenerators.lean}.

The transported local chart descends from $U_\chi.R$
to $B_\ell=U_\chi.R/P_\chi$.  Its constant coefficient
is the descended cubic basepoint.  It is graded for
$W_\ell$: matrix entries have their prescribed weights,
their inverse matrices have the same entry weights by the matrix
identity $PP^{-1}=I$, and the subtracted constant
entry is nonzero only in weight zero.
The onto-linear-coordinate and cotangent-injectivity
conditions descend along the quotient, directly from
the same conditions of the local chart.
This checks the entire argument of
\path{UnitaryThreeGroups/ExistsOriginAndChartOfIsLocPres.lean}
with the changed diagonal basepoint.

Adjoin the free variables $Y_t$ to this chart and then
restrict it to the graded subalgebra $A$.
The exact generic results
\path{CommutativeAlgebra/ChartAdd.lean}, \texttt{chart\_add},
and \path{CommutativeAlgebra/ChartGr.lean},
\texttt{chart\_gr}, apply to (G51.7), its extended local
kernel and its origin.  They preserve the specified
graded chart, constant coefficient, linear-coordinate
spanning and saturated cotangent injectivity.
The numbers of chart variables become
\[
 (1+a_{-2},\ 2+a_{-1},\ 3+a_0,\ 4+a_1,\ 2+a_2),
\]
which are exactly (G51.5).  Property 2 follows.

\proofstep{Actual entries and their homogeneous density.}
The global presentation in AC-G50 makes every actual
embedded entry homogeneous.  A homogeneous lift through
(G51.7) therefore lies in $A$.
An actual three-dimensional entry is the sum of the
two relevant entries in the first block row of the
embedding, so it too has the required homogeneous lift.
This proves the first assertion of property 3.

Write $r=r^{\rm univ}$ for the normalized representation.
Since $r(c).\mathrm{mat}=I$ and $\eta(c)=-1$, the
semidirect multiplication law gives
\[
                 (\theta r(g).\mathrm{mat})_{ij}
                                  =r(cg).\mathrm{mat}_{ij}.
                                                               \tag{G51.10}
\]
Here $\theta$ acts on the left matrix component as in
(G51.1).  Thus every embedded entry is either an actual
entry of $r(g)$, an actual entry of $r(cg)$, or zero.
Every coordinate image $q(X_s)$ or $q(Y_t)$ is consequently
an $O$-linear combination of actual entries of its weight,
with a scalar subtracted only if that weight is zero.
Include the empty product $1$ among the weight-zero
entry products.  Products of these coordinate images
then lie in the span of entry products of their total
weight.

For a homogeneous power series of weight $i$, truncate
at total degree $n$.  Its finite truncation still has
weight $i$.  All coordinate images lie in $\mathfrak m_S$,
so the remaining series maps into $\mathfrak m_S^n$.
The finite truncation has the asserted entry-product
expression.  This proves the density assertion for every
$n$, including $n=0$.
It is the full elementary proof of
\path{UnitaryThreeGroups/ExistsSpanEntMonApproxOfMemWtPart.lean},
with its only seed-specific vanishing assertion checked
directly for the diagonal cubic seed.

\proofstep{The point and dual-point fields.}
The ring $\widetilde R$, its presentation $W_A$, the
origin, the surjection $q$, the homogeneous equations
$g_u$ and the equality (G51.8) meet every algebraic
hypothesis of the two general theorems
\path{UnitaryThreeGroups/ExistsTorusTwistEqLocalPointOfWeightedPres.lean}
and
\path{UnitaryThreeGroups/ExistsDualNumberPointOfPointDerivationVanishing.lean}.
Neither statement contains a Galois representation,
residual seed, or rational-character hypothesis.
Their remaining hypothesis is exactly (G51.6), by
(G51.9).  The first theorem produces the torus translate
and its actual local map from $S$; the second produces
the actual map from $S$ to dual numbers with its prescribed
first component and derivative on $A$.
They prove property 4 with precisely the displayed
antecedent.  They do not assert (G51.6) itself.

\proofstep{The integral basepoint derivative field.}
We prove property 5, including its integral compatibility.
Write $D:S\to O/a$ for the given derivation.
Let $D_1:U_\chi.R\to O$ be the given integral local
derivation at a local cubic basepoint $x_\ell$ which
annihilates $P_\chi$.
By local framed uniqueness, any $O$-valued map
specializing the universal local representation to
the cubic basepoint is the fixed local basepoint map;
thus $x_\ell$ is the one used in the chart.
The derivation $D_1$ descends to a derivation
$D_\ell:B_\ell\to O$ at its origin.

Every coordinate of $W_A$ maps in $S$ to an embedded
entry difference.  AC-G51-diagonal shows that its
$D$-derivative vanishes unless its coordinate weight is
$-1$.  The same is true for the coordinates of
$W_\ell$ and $D_\ell$.
For each local coordinate, the unit factor in (G51.1)
is identical before and after restriction.
The assumed compatibility of $D$ and $D_1$ on actual
matrix entries therefore implies compatibility on
every local embedded coordinate:
\[
 D_\ell(X_s)\bmod a=D(q(X_s)).                               \tag{G51.11}
\]

Choose an integral value $d_s=D_\ell(X_s)$ for every
local coordinate.  For each extra coordinate of
weight $-1$, choose any lift $d_t\in O$ of its
$D$-value.  For every extra coordinate of other weight,
put $d_t=0$.
On $P$ define
\[
 L(F)=\sum_z [z]F\,d_z,                                    \tag{G51.12}
\]
where $[z]F$ is the coefficient of the degree-one
monomial $z$.  The coefficient-of-product identity
shows that $L$ is an integral derivation at constant
term.  On the local power-series subring it equals
$D_\ell$ composed with its presentation map: such
derivations are determined by their values on the
finitely many variables.  Concretely, they kill the
square of the variable ideal and all constants, so
the degree-one expression determines their value on
every power series.

It follows that $L$ kills the local presentation kernel.
It also kills the ideal it generates in $P$: the
generators have zero constant term, and the derivation
rule kills each of their products with an arbitrary
power series.  Thus $L$ descends to $\widetilde R$,
and then restricts to an integral derivation $D'$ on
$A$.  It is of weight $-1$, because in (G51.12) only
coordinates of weight $-1$ have nonzero assigned value.

Modulo $a$, its values on every coordinate agree with
$Dq$.  Both maps are derivations at the origin; the
same degree-one determination therefore gives
\[
                     D(q_A(f))=D'(f)\bmod a\qquad(f\in A).
                                                               \tag{G51.13}
\]
This proves property 5.
It is also a direct verification of every hypothesis
of the generic
\path{UnitaryThreeGroups/ExistsBasePointDerivationOfWeightHomogeneous.lean},
whose full proof implements (G51.11)--(G51.13).
Notice that no global equation was imposed on
$\widetilde R$ when lifting $D'$; precisely this relative
ambient ring makes the integral lifting statement true.

All model fields are now verified.  The proof has used
the source's general commutative-algebra theorems with
their literal hypotheses, and separately replaced the
two diagonal-seed calculations, for density and integral
derivatives, by the complete arguments above.
\end{proof}

\paragraph{Exact remaining arithmetic requirements.}
The construction closes the full graded-model existence
premise at the coefficient primes covered by AC-G49-inert.
It does not prove the numerical or geometric hypotheses
$M1$--$M5$ of \texttt{LemmaM}.
In particular the cotangent-spanning assertions in $M1$,
finite degree-zero condition $M2$, actual Type I automorphy,
and the required vanishing of the projected mixed tangent
classes still require their respective arithmetic inputs.
The actual mixed coefficient identified in AC-G45 and
AC-K29-tangent has not become a rational Dirichlet twist.
The local chart at split coefficient primes with nontrivial
local $\chi$ remains outside this theorem.

\paragraph{Geometric dictionary.}
The residual finite local systems have the finite-cover
interpretation in the already audited SGA~1, Expos\'e~V,
\S\S4 and~7.  Here the graded ambient algebra is an
explicit relative deformation construction: it keeps
the local equations and records the global equations
separately.  Its tangent lifts and chart are justified
by the displayed algebra and the exact source helpers,
not by the finite-cover dictionary alone.  No new
archival locator or motivic realization is inferred.
